% !PRINTED 63
\weberpage{63}
%<<CONT

\begin{equation}
 n\Phi(x_1,x_2,\ldots,x_m)
 =x_1\Phi'(x_1)+x_2\Phi'(x_2)+\cdots+x_m\Phi'(x_m),
\tag{5}
\end{equation}
wovon die Formel (15) des vorigen Paragraphen ein specieller Fall ist, und für $\nu=2$:
\begin{equation}
 n(n-1)\Phi(x_1,x_2,\ldots,x_m)
 =
 \sum_{i,k=1}^{m}x_i x_k\Phi''(x_i,x_k),
\tag{6}
\end{equation}
worin die Summe von $i=1$ bis $i=m$ und von $k=1$ bis $k=m$ zu erstrecken ist, so dass jedes Glied mit ungleichen $i,k$ zweimal in der Summe auftritt.

Setzen wir, wenn die $\alpha$ der Bedingung (4) unterworfen sind,
\begin{equation}
 \Phi_\nu(\xi,x)=
 \sum {\xi_1^{\alpha_1}\xi_2^{\alpha_2}\cdots\xi_m^{\alpha_m}
 \over \Pi(\alpha_1)\Pi(\alpha_2)\cdots\Pi(\alpha_m)}
 D_{\alpha_1,\alpha_2,\ldots,\alpha_m}\Phi,
\tag{7}
\end{equation}
so ist nach (4) des \S~16:
\begin{equation}
 \Phi(x+\xi)=\Phi(x)+\Phi_1(x,\xi)+\Phi_2(x,\xi)+\cdots+\Phi_n(x,\xi),
\tag{8}
\end{equation}
und da die linke Seite ungeändert bleibt, wenn $x$ mit $\xi$ vertauscht wird, so ergiebt sich die Relation:
\begin{equation}
 \Phi_{n-\nu}(x,\xi)=\Phi_\nu(\xi,x),
\tag{9}
\end{equation}
also insbesondere
\begin{equation}
 \Phi_n(x,\xi)=\Phi(\xi).
\tag{10}
\end{equation}
Die Function $\Phi_\nu(x,\xi)$ wird, als Function von $x$ betrachtet, die \gesperrt{$\nu$te Polare} der Function $\Phi$ für das Werthsystem $\xi$ genannt.

Wir wollen die Formel (3) noch für den Fall einer binären Form ($m=2$) specialisiren. Wir bezeichnen die Variablen mit $x,y$ und setzen zur Abkürzung:
\[
 \Phi(x,y)=u,\qquad D_{h,\nu-h}\Phi=u_h
\]
und erhalten aus (4):
% [source-original sic] Printed p.63 cites (4), although the displayed specialization is derived from (3); German is preserved verbatim and English gives the mathematically correct reference.
\begin{equation}
 {\Pi(n)\over\Pi(n-\nu)}u
 =
 \sum_{h=0}^{\nu}{\Pi(\nu)\over\Pi(h)\Pi(\nu-h)}
 u_h x^h y^{\nu-h},
\tag{9}% [source-original sic] Printed p.63 repeats (9) after (9) and (10); retained verbatim rather than silently renumbered.
\end{equation}
worin $\nu$ jeden beliebigen Werth, der nicht grösser als $n$ ist, annehmen kann.
